\subsection{Finite-duration quantum readout on the same scalar action}
\label{sec:source-scalar-finite-instrument}

The centered instrument admits a finite-duration realization while the field
continues to evolve. The control Hamiltonian and its precision remain supplied.
The result below bounds the resulting probabilities; it does not identify a
classical execution with a quantum outcome history.

Keep the positive scalar operator $\mathsf A$, mass inner product, original
vacuum, coherent velocity $v$ and detector $g$ used above. In canonical
coordinates $Q=M^{1/2}q$, $P=M^{-1/2}p$, write
\[
 H_0=\tfrac12(P^TP+Q^T\Omega^2Q),\qquad
 \Omega^2=M^{1/2}\mathsf A M^{-1/2}.
\]
The pointer has zero free Hamiltonian. During a rectangular pulse of duration
$\delta>0$, apply the simultaneous local coupling
\[
 H_{\rm int}=\frac1{2\delta}\sum_{i\in\operatorname{supp}g}
 Z_i\otimes m_i g_iq_i.
\]
The coupling strength scales as $1/\delta$; pulse duration alone does not
bound implementation cost or required control strength.

\begin{proposition}[Exact finite pulse on the prepared pointer code]
\label{prop:source-scalar-finite-pulse}
On the GHZ code spanned by $|0\rangle^{\otimes32}$ and
$|1\rangle^{\otimes32}$, every $Z_i$ acts as the same logical $Z$.
Define $S_\delta=\operatorname{sinc}(\delta\sqrt{\mathsf A}/2)$ and
$g_\delta=S_\delta g$, with $\operatorname{sinc}(0)=1$. The full pulse
propagator on that code is
\[
 e^{-i\delta H_0/2}\,e^{i\theta_\delta}
 e^{-iZ\otimes\Phi(g_\delta)/2}\,e^{-i\delta H_0/2},
\]
where
\[
 \theta_\delta=
 \frac{\langle g,[\delta\mathsf A^{-1}
 -\mathsf A^{-3/2}\sin(\delta\sqrt{\mathsf A})]g\rangle_M}{8\delta^2},
 \qquad 0\le\theta_\delta\le\frac{\delta\langle g,g\rangle_M}{48}.
\]
Pointer preparation and parity readout therefore realize the centered
instrument with smear $g_\delta$ at the pulse center, together with the
surrounding free field evolution.
\end{proposition}
\begin{proof}
In each logical pointer sector the field is a finite collection of harmonic
oscillators with a constant linear force of opposite sign. Its exact Weyl
solution has the averaged interaction-picture linear observable and a central
second-order phase. Averaging symmetrically about the pulse center cancels
the momentum term and multiplies each position mode by
$\operatorname{sinc}(\delta\omega/2)$. The central phase is quadratic in
the force, so both logical sectors have the displayed same phase. These
identities can be computed on the oscillator Schwartz core and extended
through the bounded unitary operators. The inequalities
$0\le x-\sin x\le x^3/6$ for $x\ge0$ give the phase bound.
The earlier parity calculation then applies with $g_\delta$.
\end{proof}

The phase is scalar on the ideal GHZ code only. On the full pointer Hilbert
space the second-order terms contain $Z_iZ_k$; they are not generally scalar.
Noise comparisons below use the full pulse channel and do not assume that a
noisy pointer stays in the code. Also, the local coupling support and the
effective equal-time smear are different: $S_\delta g$ generally extends
outside the original detector region. The instantaneous regional-algebra
statement for $g$ cannot be transferred to this averaged smear.

For centers $t_1<\cdots<t_{21}$, the interaction-picture observable is
$B_j=\Phi(g_\delta)(t_j)$. The same two Kraus formulas define a joint
field/record instrument whose classical record prefixes are consistent.
Its unconditional sine response is
\[
 R_j^\delta=\frac12e^{-\nu_\delta/2}\sin\mu_\delta(t_j)
 \prod_{k<j}\cos\frac{c_\delta(t_j-t_k)}2,
\]
where
\begin{align*}
 \mu_\delta(t)&=\langle g,S_\delta\mathsf A^{-1/2}
               \sin(t\sqrt{\mathsf A})v\rangle_M,\\
 \nu_\delta&=\tfrac12\langle g,S_\delta^2\mathsf A^{-1/2}g\rangle_M,\\
 c_\delta(t)&=\langle g,S_\delta^2\mathsf A^{-1/2}
               \sin(t\sqrt{\mathsf A})g\rangle_M.
\end{align*}
The centered ideal baseline remains $1/2$. The ideal nonselective pulse
preserves canonical first moments relative to free evolution and adds
$\|g_\delta\|_M^2/8$ to the expected field energy. This is not a bound on
pointer, controller or noisy implementation energy.

\begin{proposition}[Duration and channel-error comparison]
\label{prop:source-scalar-finite-noise}
Put $G_0=\langle g,g\rangle_M$, $G_1=\langle g,\mathsf A g\rangle_M$
and $V_0=\langle v,v\rangle_M$. The finite pulse and the instantaneous
instrument on the same continuously evolving field satisfy
\begin{align*}
 |\mu_\delta-\mu_0|&\le\delta^2\sqrt{G_1V_0}/24,\\
 |\nu_\delta-\nu_0|&\le\delta^2\sqrt{G_0G_1}/24,\\
 |c_\delta(d)-c_0(d)|&\le\delta^2G_1|d|/12,\\
 |\cos(c_\delta(d)/2)-\cos(c_0(d)/2)|
 &\le\delta^2G_0G_1d^2/48.
\end{align*}
Suppose $t_k\in[a_k,b_k]$ are the inherited ordered center intervals and
$u_j$ bounds the unmeasured continuous response magnitude there. Set
\[
 B_\delta=\delta^2\left(\frac{\sqrt{G_1V_0}}{48}
                  +\frac{\sqrt{G_0G_1}}{96}\right),\qquad
 A_j=\frac{\delta^2G_0G_1}{48}\sum_{k<j}(b_j-a_k)^2.
\]
Then $|R_j^\delta-R_j^0|\le B_\delta+u_jA_j$.

Assume each initial state has trace distance at most $\epsilon_0$ from its
specified preparation and each of the 129 declared fault locations per
readout has channel distance at most $\epsilon$ in half the diamond norm.
The error of one outcome probability through readout $j$ is at most
$\epsilon_0+129j\epsilon$. For the difference between noisy intervention
and noisy baseline probabilities, add $2\epsilon_0+258j\epsilon$.
\end{proposition}
\begin{proof}
Use $|\operatorname{sinc}x|\le1$ and
$|1-\operatorname{sinc}x|\le x^2/6$ in the finite spectral calculus.
Cauchy--Schwarz gives the mean and variance bounds; in particular
$\langle g,\sqrt{\mathsf A}g\rangle_M\le\sqrt{G_0G_1}$.
The inequality $|\sin(d\omega)|\le|d\omega|$ gives the commutator bound.
Both commutators have absolute value at most $G_0|d|$, and
$|\cos x-\cos y|\le|x-y|(|x|+|y|)/2$ proves the last line.
The unmeasured probability is half a damped sine. Its mean contribution is
bounded by half the mean error and its variance contribution by one quarter
of the variance error, giving $B_\delta$. Telescope the attenuation
products, whose factors have magnitude at most one, to obtain $A_j$.
Finally telescope the implemented channels and use trace-distance
contractivity, retaining the complete field, pointers and classical records.
The two preparations may have different errors; hence both probability
errors enter their difference. A noisy baseline is not fixed at $1/2$.
\end{proof}

The channel-error assumptions do not follow from a bound on coefficients of
an unbounded field Hamiltonian. The declared model places the 32 local pulse
fault channels at its common endpoint. Pointer-only operations commute with
the free field flow, and later operations preserve retained records.
These are explicit control and record assumptions. Half-diamond closeness
controls bounded probabilities, not energy or noisy first moments.

\paragraph{A nonzero duration and noise window.}
The independently replayed certificate retains all 21 centers and all 210
earlier-readout pairs. It uses a common pulse duration
$1/1000\le\delta\le1/100$, pointer-only operation duration
$\gamma=1/10000$, and bounds
$\epsilon\le10^{-7}$, $\epsilon_0\le10^{-5}$.
Each readout has 64 sequential preparation operations, 32 simultaneous
local field couplings and 33 read/decode operations. Its occupied window
is $\delta+97\gamma\le0.0197$, with all 2,709 operation locations retained.
Every choice of centers in the inherited intervals leaves a gap greater
than $0.04417$ between these occupied windows.

Combining the pulse and noise bounds with the earlier same-instrument
clock comparison retains all fifteen resolved responses. The smallest
retained magnitude exceeds $0.0003659$; at the sixteenth readout the noisy
intervention-minus-baseline response is negative with magnitude exceeding
$0.0082353$. Its total comparison error is below $0.0012705$.
The final parity record becomes available in
$[1.349935903357,1.349945669638]$ at the maximum pulse duration, after its
center-time interval. The ideal field-energy increase per pulse is enclosed
by $[0.007652475342308,0.007656300183578]$.

The package \texttt{code/source\_scalar\_finite\_instrument/} contains the
producer, independent verifier and mutation controls. It checks the complete
parent action and instrument, moments, chronology, fault counts and outward
error bounds. The observer-like detector has local pointers, coupling ports,
retained readback records, fresh-pointer resets and declared control. The
certificate is an operator and error calculation, not a sampled run.

Initial vacuum/coherent preparations are supplied at model time zero; the
sixteen earlier ideal preparation factors are not given a finite-duration
implementation by this result. The inherited clock is a conditional
reference-mean enclosure, not a noisy or branch-conditioned clock read from
pointer bits. Quantum preparation, entangling operations, Born readout,
physical regions and time units remain supplied or unattached. No native
seam implementation, spatial continuum or interacting-field limit follows.
